% ---------------------------------------------------------------------------
% Capítulo de considerações finais.
% Base: capítulos de introdução, metodologia, resultados, requisitos,
% desenvolvimento e avaliação (situação de 17/11/2026). Os números repetem os
% dos capítulos; não foram recalculados. A única fonte externa é a Portaria
% SGD/MGI nº 7.076/2024 (docs/referencias-verificadas.md, seção
% "Considerações finais (17/11/2026)").
% ---------------------------------------------------------------------------
\begingroup
\setlength{\emergencystretch}{3em}

\chapter{Considerações finais}
\label{cap:consideracoes}

Este capítulo retoma os objetivos do trabalho, resume os achados das duas
revisões, as contribuições e as limitações, e indica os trabalhos futuros.

% ===========================================================================
\section{Objetivos alcançados}
\label{sec:consideracoes-objetivos}

O objetivo geral era desenvolver um protótipo de buscador OSINT para a
polícia judiciária apurar fraude em contratações públicas brasileiras, com
requisitos fundamentados numa revisão sistemática, na legislação e no
levantamento com especialista, e com controles de preservação da prova desde
a coleta (Seção~\ref{sec:intro-objetivos}). O protótipo foi desenvolvido
(Capítulo~\ref{cap:desenvolvimento}). Cada história do \textit{backlog} aponta
o requisito que atende e a evidência que a sustenta
(Seção~\ref{sec:requisitos-historias}), e toda captura entra na cadeia de
custódia (Seção~\ref{sec:desenv-arquitetura}). Os objetivos específicos foram
atendidos assim:
\begin{enumerate}
	\item \textbf{RSL 1.} A revisão incluiu 57 estudos primários, 23 da busca
	nas bases e 34 do \textit{snowballing} (Seção~\ref{sec:rsl1-selecao}), e
	sintetizou as ferramentas, as técnicas de IA, os requisitos legais e
	éticos, as ameaças e a avaliação
	(Seções~\ref{sec:rsl1-qp1-qp2} a~\ref{sec:rsl1-qp5}). Os resultados são
	provisórios, pelas pendências descritas na
	Seção~\ref{sec:consideracoes-limitacoes}.
	\item \textbf{RSL 2.} A revisão sintetizou 129 estudos, em 130 relatos
	(Seção~\ref{sec:resultados-rsl2-selecao}), nas dimensões de
	produtividade, qualidade, segurança, custo e processo
	(Seções~\ref{sec:resultados-rsl2-direcao}
	a~\ref{sec:resultados-rsl2-qp5}), e terminou em dez práticas candidatas
	para o desenvolvimento (Seção~\ref{sec:resultados-rsl2-praticas}). Os
	resultados também são provisórios.
	\item \textbf{Requisitos e \textit{backlog}.} A síntese da RSL 1 reuniu 27
	requisitos candidatos. Dezoito já tinham correspondente no levantamento,
	e seis dos outros nove viraram requisitos novos
	(Seção~\ref{sec:requisitos-novos}). Os requisitos foram escritos como 39
	histórias de usuário, em sete épicos, e priorizados pelo MoSCoW
	(Seções~\ref{sec:requisitos-historias}
	e~\ref{sec:requisitos-priorizacao}).
	\item \textbf{Comparação com ferramentas.} O buscador foi comparado com
	sete ferramentas, pela documentação oficial de cada uma
	(Seção~\ref{sec:requisitos-mercado}). A comparação mostrou o que falta a
	elas no domínio deste trabalho e o que falta ao protótipo, organizado em
	três horizontes do \textit{backlog}.
	\item \textbf{Protótipo.} Das 39 histórias, 25 estão feitas, 6 em parte, 3
	por fazer e 5 ficaram fora desta entrega. As 11 histórias \textit{Must
	have} estão feitas (Seção~\ref{sec:requisitos-priorizacao}). O protótipo
	tem coleta só passiva, cadeia de custódia com SHA-256 em cada captura,
	trilha de auditoria encadeada, dez tipologias por regras explícitas, com
	alertas tratados como pistas, e extração por IA local, desligada por
	padrão (Capítulo~\ref{cap:desenvolvimento}). As ameaças tratadas estão
	ligadas a testes automatizados (Seção~\ref{sec:desenv-ameacas}), e os
	conectores foram conferidos contra as fontes reais
	(Seção~\ref{sec:desenv-verificacao}).
	\item \textbf{Avaliação.} O objetivo era planejar a avaliação com
	usuários policiais, e o protocolo está pronto
	(Seção~\ref{sec:avaliacao-protocolo}). A avaliação não foi executada. Ela
	depende da validação do orientador e da decisão do sistema CEP/CONEP
	(Seção~\ref{sec:avaliacao-etica}).
\end{enumerate}
% PENDENTE: se a avaliação com usuários for executada, atualizar o item 6 e
% o objetivo específico 6 da introdução.

% ===========================================================================
\section{Principais achados das revisões}
\label{sec:consideracoes-achados}

\textbf{RSL 1.} Nenhum dos 57 estudos é do Brasil ou de outro país da
América Latina, e nenhum usa dados em português
(Seções~\ref{sec:rsl1-caracteristicas} e~\ref{sec:rsl1-qp5}). Nenhum trata de
fraude ou corrupção em contratações públicas
(Seção~\ref{sec:requisitos-mercado}). A avaliação mede o modelo, e quase
nunca o trabalho do investigador: só 7 estudos avaliam com usuários. A
preservação da evidência também é exceção: só 7 estudos tratam
explicitamente da cadeia de custódia, e nenhum avalia se a evidência é
aceita em juízo (Seção~\ref{sec:rsl1-qp5}). Os estudos tratam a saída da IA
como pista, com a decisão e a explicação com o humano, e descrevem um modelo
de ameaças com três alvos: o investigador, a ferramenta e os dados
(Seção~\ref{sec:rsl1-qp3-qp4}). Nenhum estudo relata teste de segurança da
própria ferramenta. A revisão confirma o núcleo probatório do levantamento,
a cadeia de custódia com a trilha de auditoria, e não dá respaldo às
tipologias de corrupção nem às fontes brasileiras, que continuam no projeto
pelo levantamento e pela legislação (Seção~\ref{sec:rsl1-rastreabilidade}).

\textbf{RSL 2.} O ganho de produtividade existe, mas é condicional e menor do
que se percebe. Os resultados negativos e nulos aparecem sobretudo quando há
medida objetiva (Seções~\ref{sec:resultados-rsl2-direcao}
e~\ref{sec:resultados-rsl2-qp1}). Na qualidade, a direção mista predomina, e
a qualidade depende da verificação humana, que a IA enfraquece
(Seção~\ref{sec:resultados-rsl2-qp2}). O código gerado não é seguro por
padrão, e quem usa IA superestima a segurança do próprio código
(Seção~\ref{sec:resultados-rsl2-qp3-qp4}). No processo, o portão humano
permanece, a verificação vira o gargalo, e as práticas de governança são
propostas ou descritas, mas nenhum estudo mede o efeito delas. Também não há
estudo sobre o desenvolvimento individual de um produto completo em contexto
acadêmico (Seção~\ref{sec:resultados-rsl2-qp5}).

% ===========================================================================
\section{Contribuições}
\label{sec:consideracoes-contribuicoes}

As contribuições são as da Seção~\ref{sec:intro-contribuicoes}: as duas
revisões sistemáticas, com a rastreabilidade dos achados para os requisitos
e para o processo; o \textit{backlog} priorizado e a comparação com
ferramentas existentes; o protótipo de código aberto, sob a licença MIT, com
o modelo de ameaças ligado aos testes; e o registro do desenvolvimento com
agentes de IA. O protótipo responde, em parte, a lacunas que a RSL 1 apontou:
usa fontes brasileiras e dados em português, une cadeia de custódia, trilha
de auditoria e grafo de vínculos, e trata o alerta e a saída da IA como
pistas que o investigador confirma ou descarta. Se isso ajuda o trabalho do
investigador ainda não foi avaliado.

% ===========================================================================
\section{Limitações}
\label{sec:consideracoes-limitacoes}

\textbf{Método das revisões.} Na etapa 1, R1 e R2 triaram de forma
independente (Seção~\ref{sec:metodologia-selecao}). Da etapa 2 em diante, a
leitura do texto completo, o \textit{snowballing}, a avaliação de qualidade,
a extração e a síntese foram feitos pela IA como terceira opinião, com regras
aceitas pelo autor (R1). O orientador (R2) ainda confirma essas regras, e a
avaliação independente de R1 e R2 nessas etapas, com o kappa de cada uma,
ainda não foi feita (Seções~\ref{sec:resultados-rsl1}
e~\ref{sec:resultados-rsl2}). Ficaram relatos sem leitura: 21 na busca e 43
no \textit{snowballing} da RSL 1, e 12 e 35 na RSL 2, não foram recuperados.
Os estudos secundários das duas revisões ainda não foram lidos. Há ainda
pares de relatos que podem ser do mesmo estudo, a decidir por R1 e R2
(Seções~\ref{sec:rsl1-caracteristicas}
e~\ref{sec:resultados-rsl2-selecao}). O recorte da etapa 2 deixou de fora,
na RSL 1, os estudos conceituais, legais e éticos sem avaliação de
ferramenta (Seção~\ref{sec:rsl1-qp3-qp4}).
% PENDENTE: R2 confirma as regras de aplicação das etapas 2 em diante; R1 e
% R2 fazem a avaliação independente e calculam o kappa. Atualizar este
% parágrafo quando isso acontecer.

\textbf{Protótipo.} Os testes automatizados e a demonstração usam só dados
sintéticos. As fontes reais passaram por um teste de fumaça, e não por uma
investigação (Seção~\ref{sec:desenv-verificacao}). A base completa do CNPJ
só entra por recorte, e não há OCR para editais escaneados. A extração por
IA foi testada com três documentos sintéticos, o que não mede a sua
qualidade. A importação do grafo no Maltego e no i2 não foi testada nos
programas. A troca da chave de criptografia não recifra as capturas antigas.
Na última rodada, os testes específicos de PostgreSQL não rodaram. Estão em
tratamento o CPF completo nas capturas brutas e na comparação entre capturas
e o CNPJ alfanumérico (Seção~\ref{sec:desenv-limitacoes}).
% PENDENTE: atualizar quando a correção do CPF/CNPJ alfanumérico terminar.

\textbf{Avaliação.} O protótipo ainda não foi avaliado com policiais
(Capítulo~\ref{cap:avaliacao}). Sem essa avaliação, não há como afirmar que
ele apoia o trabalho do investigador nem que se aproxima das ferramentas
comparadas (Seção~\ref{sec:requisitos-mercado}).

% ===========================================================================
\section{Trabalhos futuros}
\label{sec:consideracoes-futuros}

\textbf{Avaliação com usuários.} O primeiro trabalho é executar o protocolo
do Capítulo~\ref{cap:avaliacao}, depois da decisão do sistema CEP/CONEP
(Seção~\ref{sec:avaliacao-etica}).

\textbf{Revisões.} Concluir a confirmação das regras por R2, a avaliação
independente de R1 e R2 com o kappa, a leitura dos estudos secundários e os
pedidos dos relatos não recuperados, e recalcular as tabelas por
\textit{script} se algum resultado mudar.

\textbf{Uso em produção.} Um órgão que queira usar o protótipo precisa de
mais do que este trabalho entrega. As medidas a seguir são decisão do autor:
\begin{itemize}
	\item recifrar as capturas antigas quando a chave de criptografia for
	trocada;
	\item fazer um teste de invasão independente, já que nenhum estudo da
	RSL 1 relata teste de segurança da ferramenta
	(Seção~\ref{sec:rsl1-qp3-qp4}) e a auditoria do protótipo foi assistida
	por IA (Seção~\ref{sec:desenv-ameacas});
	\item elaborar, com o órgão, o relatório de impacto à proteção de dados
	pessoais e a análise da LGPD para o tratamento feito pela ferramenta;
	\item avaliar o carimbo do tempo emitido no âmbito da ICP-Brasil para as
	capturas e a trilha, em lugar do relógio do servidor;
	\item acrescentar a entrega contínua à integração contínua já
	configurada nos dois repositórios (Seção~\ref{sec:desenv-verificacao});
	\item integrar o acesso à Conta gov.br. Pela Portaria SGD/MGI
	nº~7.076/2024, quem pede essa integração é o gestor público do órgão, e
	o sistema precisa estar hospedado num domínio oficial de governo
	\cite[art.~3º e art.~4º, § 1º]{brasil2024portariasgd7076}. Por isso, essa integração só
	pode ser pedida pelo órgão que adotar a ferramenta.
\end{itemize}

\textbf{Fontes e tipologias.} Acrescentar as demais fontes do levantamento,
como diários oficiais, tribunais de contas, juntas comerciais e cartórios,
e as tipologias ainda não implementadas (Tabela~\ref{tab:backlog}). Nas
fontes já ligadas, falta baixar o edital pelo próprio conector, ler os
editais escaneados por OCR, buscar no CNPJ as outras empresas de um sócio do
caso e ler arquivos do TSE acima do teto atual
(Seção~\ref{sec:desenv-limitacoes}).

\textbf{IA.} Medir a qualidade da extração por IA com um conjunto de editais
anotado, que este trabalho não tem (Seção~\ref{sec:desenv-limitacoes}).

\endgroup
